\section{P21: three non-equivalent connectivity questions}\label{sec:p21}

The uniform theorem concerns one selected one-copy margin fiber.  The P21
example is included to show what that fiber cannot certify.  It distinguishes:
\begin{enumerate}
  \item connectivity of one fixed nonnegative fiber;
  \item generation of the complete integer relation lattice;
  \item connectivity of every nonnegative fiber, equivalently the Markov
        basis requirement.
\end{enumerate}
The first fixes one right-hand side.  The second allows signed intermediate
combinations.  The third varies the right-hand side and must preserve
nonnegativity throughout \cite{DiaconisSturmfels1998}.  Algebraic
formulations of grid-tiling flip ideals also precede this example
\cite{Chin2019,GrossYamzon2021}; the claims below are only for the frozen
72-variable configuration.

\subsection{Frozen configuration}

The source-locked specimen is \emph{4x4 square target}, signed LS 11-02, with
SHA-256
\[
\mathtt{2248f6e9d0d4cb793ef2ff8069a37ae62ed67629cee6263a0f7071ae4a274120}.
\]
Its target is the \(5\times5\) square.  We freeze the zero-based piece-index
row
\[
 E=(4,5,12,13),
\]
whose areas are \((7,6,6,6)\).  Only placements participating in the selected
margin fiber are retained.  Their block counts are \(8,24,16,24\), for 72
variables.

For each participating placement \(P\) in named block \(i\), the
configuration matrix \(A\) has column
\[
 e_i\mid X_{\rm row}(P)\mid X_{\rm col}(P).
\]
Thus \(A\) is a \(14\times72\) integer matrix of rank 12 and
\[
 K=\ker_{\ZZ}(A)
\]
has rank 60.

For \(s=1,2,3,4\), let \(M_{\leq s}\) be the integer span of all squarefree
one-copy relations whose two placement tuples differ in at most \(s\) named
piece positions.  This \emph{piece support} is not toric degree, Graver
support, or a claim of Graver primitiveness.

\subsection{The fixed 32-node fiber}

The chosen fiber contains 32 nodes: 8 exact tilings and 24 non-tiling margin
matches.  Its complete support filtration is:

\begin{center}
\small
\begin{tabular}{@{}lrrrrl@{}}
\toprule
Family & New rels. & Lattice rank & \(\rank K/M\) & Components & Verdict\\
\midrule
\(M_{\leq2}\) & 144 & 34 & 26 & 16 & 16 trapped nodes\\
\(M_{\leq3}\) & 932 & 51 & 9 & 4 & all reach tilings\\
\(M_{\leq3}+\langle O_{714}\rangle\) & one orbit & 54 & 6 & 1 & connected\\
\(M_{\leq4}\) & 7,444 & 60 & 0 & 1 & complete graph\\
\bottomrule
\end{tabular}
\end{center}

The relation counts are new at the indicated exact support.  The complete
\(M_{\leq3}\) and \(M_{\leq4}\) catalogues contain 1,076 and 8,520 relations,
respectively.  At support at most two there are sixteen two-node components:
eight pair a tiling with an extra, while eight contain the sixteen trapped
nodes.  At support three there are four eight-node components, but every node
reaches a tiling.  At support four the graph is complete.

\begin{proposition}[One connected fiber does not generate the lattice]
\label{prop:p21-one}
The family \(M_{\leq3}+\langle O_{714}\rangle\) connects the fixed 32-node
fiber, but has lattice rank \(54<60=\rank K\).
\end{proposition}

\begin{proof}
The order-eight target action partitions the 7,444 support-four relations
into 1,119 complete orbits.  Exactly 36 single orbit modules connect the four
components of the support-three graph.  Direct Hermite/Smith reduction gives
rank 54, with torsion-free quotient rank six, for every one of these choices;
\(O_{714}\) is the canonical stored witness.
\end{proof}

\subsection{The missing integral channels}

The quotient \(Q=K/M_{\leq3}\) is torsion-free of rank nine.  Among the 1,119
support-four orbits, 87 cyclic orbit modules have quotient rank three.  They
fall into three exact integral families of sizes
\[
 6,\quad3,\quad78.
\]
Exactly three complete target-action orbits are necessary and sufficient to
generate \(Q\) integrally.  Every minimum triple chooses one orbit from each
family, giving \(6\cdot3\cdot78=1,404\) minimum triples.  The stored canonical
triple is \((283,603,714)\), with determinant \(-1\) in the quotient basis.

The fixed 32-node fiber activates \(0,0,36\) orbits in these three families.
It is blind to the first two necessary integral channels.  These numbers are
configuration-specific finite facts, not a general equivariant-basis
theorem; compare the general lifting context in
\cite{RauhSullivant2019}.

\subsection{Lattice generation is not Markov connectivity}

All 60 Smith factors for \(M_{\leq4}\subseteq K\) are one, hence
\[
 M_{\leq4}=K
\]
integrally.  This permits signed decompositions of every relation.  It does
not imply that a decomposition can be ordered while staying nonnegative.

\begin{proposition}[All-right-hand-side no-go]\label{prop:p21-markov}
Although \(M_{\leq4}=K\), the squarefree one-copy support-at-most-four family
is not a Markov basis for the frozen 72-variable matrix.
\end{proposition}

\begin{proof}
The complete quadratic census has 218 indispensable binomials: 144
cross-block quadrics occur in the one-copy family, while 74 within-block
quadrics do not.  A canonical missing quadratic lies in a two-monomial fiber
with block count \((2,0,0,0)\).  No one-copy cross-block move applies in that
fiber, so it is disconnected.  Indispensability is certified using the
standard binomial-fiber criterion
\cite{AokiTakemuraYoshida2008,CharalambousThomaVladoiu2016}.
\end{proof}

A separate two-monomial fiber with block count \((0,0,0,7)\) forces an
indispensable binomial of toric degree seven.  Therefore every Markov basis
has degree \emph{at least} seven.  The bounded minimum-generator census is:

\begin{center}
\begin{tabular}{@{}rrr@{}}
\toprule
Toric degree & Betti fibers & Minimum generators forced\\
\midrule
2 & 218 & 218\\
3 & 996 & 996\\
4 & 5,705 & 5,727\\
\bottomrule
\end{tabular}
\end{center}

At least 6,941 minimal generators are forced through degree four.  No upper
bound of seven, exact Markov degree, complete basis, Graver basis, universal
Markov basis, or blockwise lifting theorem follows.
